\begin{subsection}{Quotient Lattices} \label{quotientlatticesect}
    In the algebro-geometric setting, we defined coordinate rings for a variety  \( V \)
    via \( \mathbb{K}[X]/J \), where \( J = I(V) \) was the vanishing ideal.
    Moreover, for each Zariski open subset \( U \subset V \) we associated
    the ring of regular functions and at a specific point \( p \in V \), the 
    local ring \( \mathcal{O}_{p} \) and the residue field \( \resf{p} \).
    The natural question is if there are similar concepts for the lattice \( \mathcal{R} \)
    and acceptance sets of members. As it turns out there is a less unique choice 
    in this setting, because both ideals and filters allow for similar constructions. 
    We will cover both here and later state in Remark \ref{dualfilterquotient}, how the constructions 
    hold for both of them.

    Let \( I \) be an ideal in a distributive lattice \( (L, \vee, \wedge) \),
    and define the relation \( a \sim b \iff \exists i \in I,\ a \vee i = b \vee i\).
    Naturally for any \( a \in L \) and \( i \in I \), \( a \vee i = a \vee i \),
    so \( (\cdot \sim \cdot) \) is reflexive. By commutativity of \( \vee \) the relation is also symmetric
    and if \( a \sim b, b \sim c \) with corresponding \( i_a, i_c \) one derives
    \begin{align*}
        a \vee i_a &= b \vee i_a \\
        i_c \vee (a \vee i_a) &= i_c \vee (b \vee i_a) \\
        a \vee (i_a \vee i_c) &= c \vee (i_c \vee i_a)
    \end{align*}
    by associativity and commutativity of \( \wedge \), hence \( a \sim c \) with \( i_c \vee i_a \)
    which is in \( I \) as \( I \) is an ideal.
    We conclude that \( (\cdot \sim \cdot) \) is an equivalence relation.
    Moreover, let \( a_1 \sim b_1 \), \( a_2 \sim b_2 \) via \( i_1, i_2 \in I \)
    then 
    \begin{align*}
        (i_1 \vee a_1) \vee (i_2 \vee a_2) = (i_1 \vee b_1) \vee (i_2 \vee b_2) \\
        (a_1 \vee a_2) \vee (i_1 \vee i_2) = (b_1 \vee b_2) \vee (i_1 \vee i_2)
    \end{align*}
    so \( (a_1 \vee a_2) \sim (b_1 \vee b_2) \). Similarly
    \begin{align*}
        (i_1 \vee i_2) \vee \left( (i_1 \vee a_1) \wedge (i_2 \vee a_2) \right)  \\
     \left ( (i_1 \vee i_2) \vee (i_1 \vee a_1) \right) \wedge \left( (i_1 \vee i_2) \vee (i_2 \vee a_2) \right)  \\
     \left ( (i_1 \vee i_2) \vee a_1) \right) \wedge \left( (i_1 \vee i_2) \vee a_2) \right) \\
     (i_1 \vee i_2) \vee (a_1 \wedge a_2)
    \end{align*}
    and doing the same with \( (i_1 \vee b_1) \wedge (i_2 \vee b_2) = (i_1 \vee a_1) \wedge (i_2 \vee a_2) \), gives 
    with \( i_1 \vee i_2 \in I  \) (since it is an ideal) that \( (a_1 \wedge a_2) \sim (b_1 \wedge b_2) \), making \( (\cdot \sim \cdot) \) 
    a compatible relation with the lattice operations and by the quotient theorem for universal algebras,
    we have a well defined "quotient lattice" \( L / I \). On a side note, if \( \vee, \wedge \) are preserved
    so are any inequalities between elements in the lattice under the equivalence relation.
    \begin{definition} \label{defquotlat}
        For an ideal \( I \) of a distributive lattice \( (L, \vee, \wedge) \)
        we call the lattice \( (L / I, \vee, \wedge) \) the quotient lattice of \( L \)
        by \( I \) (or just quotient lattice if it is clear from context).
    \end{definition}

    What does that mean in the case of risk measures?
    By Proposition \ref{acceptancesetareideals} any acceptance set induces an ideal \( \acceptidealb{A} \).
    Clearly all elements in \( \acceptidealb{A} \) are equivalent to each other in \( \extrm / \acceptidealb{A} \),
    and since \( \rho_\bot \in \acceptidealb{A} \) it is sensible to summarize them under this representant \( [ \rho_\bot ] \).
    All elements \( \rho \in \extrm \) which dominate members of \( \eta \in \acceptidealb{A} \) in the sense \( \eta \leq \rho \)
    satisfy \( \rho \vee \eta = \rho \), hence they are in distinct classes.
    The most interesting phenomenon happens for those \( \rho \in \extrm \), which neither satisfy \( \rho \leq \eta \)
    nor \( \rho \geq \eta \). For example if two risk measures \( \rho_1, \rho_2 \) are different on some subset \( S \subset \mathcal{X} \)
    but otherwise coincide, they could become equivalent if there's some \( \eta \in \acceptidealb{A} \), which dominates both of them on \( S \).
    Additionally since they are not in the ideal, we know that there are \( X, Y \in A \), s.t. \( \rho_1(X) > 0, \rho_2(Y) > 0 \),
    but \( \eta(X) \leq 0, \eta(Y) \leq 0 \), so if \( \rho_1 \sim \rho_2 \) then they must coincide for all \( X \in A \) which they
    reject, thus we could still evaluate equivalence classes of risk measures \( \rho \) with respect to their rejection sets.
    %A similar observation holds also for \( \extrm / \rejectidb{A} \) and the equi
    \newline
    For concrete computations this is somewhat inadequate, since we analyze risk measures both with respect to their values they accept and those
    they don't. In general just because two risk measures \( \rho_1, \rho_2 \in \extrm \) agree on their rejection set when restricted \( A \), 
    does not imply a globally the existence of some \( \eta \in \acceptidealb{A} \) with \( \eta \vee \rho_1 = \eta \vee \rho_2 \). 
    The most valuable information which therefore can still be extracted from this setting is given less by concrete values and more by families
    of risk measures given either by ideals and filters.
    \newline
    \newline
    The behavior of ideals from our larger lattice \( \extrm \) should then translate into ideals of \( \extrm / \acceptidealb{A} \), similarly
    how ideals from a general ring translate into ideals of a quotient ring with respect to some ideal.
    Also we would like to have a relation between \( \text{Spec } \extrm \) and \( \text{Spec } \extrm / \acceptidealb{A}\).
    Fortunately, this is a pattern which is shared by distributive lattices. 
    \begin{definition}
        Let \( L \) be a lattice and denote by \( \idl{L} \) the set of ideals in \( L \)
        and by \( \flt{F} \) the set of filters of \( L \).
    \end{definition}
    \begin{proposition}
        \( \idl{L} \) is itself a distributive lattice under the operations
        \begin{align*}
            I \wedge J = \{ i \wedge j | i \in I, j \in J \} \\
            I \vee J = \{ i \vee j | i \in I, j \in J \}
        \end{align*}
        for \( I, J \in \idl{L} \). Moreover, the meet simplifies to setwise intersection
        \begin{align*}
            I \cap J = I \wedge J
        \end{align*}
        and if \( L \) was bounded then \( \idl{L} \) is complete.
    \end{proposition}
    \begin{proof}
        See Lemma 2.5.1, Corollary 1.3.2, Corollary 1.3.15 in \cite{genlattheorygraetzer}.
    \end{proof}
    \begin{lemma} \label{lemmaprincipalequivalence}
        Let \( (i)^\downarrow \subset L \) be a principal ideal with generator \( i \) in a distributive lattice \( L \).
        Then for \( a, b \in L \) it holds
        \[
            [a] = [b] \iff a \vee i = b \vee i
        \]
    \end{lemma}
    \begin{proof}
        Assuming the RHS implies the LHS by Definition. For the converse, let \( j \in (i)^\downarrow \),
        s.t. \( a \vee j = b \vee j \), applying the join with \( i \) to both sides then reduces to the statement.
    \end{proof}
    \begin{proposition} \label{correspondencetheorem}
        For distributive lattices the correspondence theorem holds, that is if \( I \in \idl{L} \)
        and \( (I)^\uparrow = \{ J \in \idl{L} | I \subseteq J \} \), then we have a bijection
        \[
            \Phi : (I)^\uparrow \mapsto \idl{L/I},
        \]
        which preserves inclusion, primeness and maximality.
    \end{proposition}
    \begin{proof}
        Being bijective is a restatement of Lemma 2.6.11 in \cite{genlattheorygraetzer} (second isomorphism theorem for universal algebras) by noting that one can substitute 
        for congruence relation the equivalence relation defined by ideals for their quotient lattices.
        As inclusion is preserved under maps, it is in particular by the projection \( \phi : L \to L / I \)
        and then \( \Psi \) being bijective maintains the inclusion of the ideals.
        If inclusions are preserved so is the property of being maximal.
        Further it holds \( \pi(J) = \pi(J)^\downarrow \) (the generated ideal by \(J\) in the quotient), because 
        if \( [a] \leq [j] \) for some \( j \in \pi(J) \), we can write it as \( [a] = [j \wedge b] \) for some \( b \in L \)
        but \( J \) is an ideal so \( j \wedge b \in J \) and then we have \( \pi^{-1}([a]) \in J \).
        Let \( P \in (I)^\uparrow \) be prime, then if \( [a \wedge b] \in \pi(P) \) also \( a \wedge b \in \pi^{-1}([a \wedge b]) \)
        is also in \( P \) and by assumption then \( a \) or \( b \) too.
    \end{proof}
    \begin{proposition} \label{quotientprojectionlattice}
        For a distributive lattice \( L \) and \( I, J \in \idl{L} \) s.t. \( I \subseteq J \), there is a well-defined lattice homomorphism
        \[
            \varphi_{J,I} : L / I  \mapsto L / J, \quad [x]_{I} \mapsto [x]_{J},
        \]
        for \( x \in L \).
    \end{proposition}
    \begin{proof}
        Let \( a,b \in L \) and \( [a]_{I}, [b]_{I} \in L / I \) be their projections into \( L / I \) and assume that \( [a]_{I} = [b]_{I} \),
        by Definition there's some \( i \in I \) s.t. \( a \vee i = b \vee i \), since \( I \subseteq J \), we also have that \( [a]_{J} = [b]_{J} \)
        making \( \varphi \) well-defined as a map. That it is also a homomorphism is a standard verification.
    \end{proof}
    Essentially, if \( I \subseteq J \) the equivalence relation of \( J \) is broader than the one of \( I \). Note that 
    \( \acceptidealb{U} \) was inclusion reversing (Corollary \ref{acceptidealincrev}) , so for \( U \subseteq V \) 
    the ideals satisfy \( \acceptidealb{U} \subseteq \acceptidealb{V} \). The intuition is that accepting \( V \) is less 
    restrictive then it is to accept \( U \), so the ideal is larger. This motivates the next observation, which should remind the reader 
    of something.
    \begin{lemma} \label{triangleidealcommutes}
        Let \( I \subseteq J \subseteq K \) be ideals in a distributive lattice \( L \), then the following 
        diagram commutes:
        \[
        \begin{tikzcd}[row sep=large, column sep=large]
            L / I \arrow[dr, "\varphi_{J, I}"] \arrow[rr, "\varphi_{K, I}"] & & L / K \\
                                                                            & L / J \arrow[ur, "\varphi_{K, J}"]  
        \end{tikzcd}
        \]
    \end{lemma}
    \begin{proof}
        By Proposition \ref{quotientprojectionlattice} the projections are all well defined
        and we need to show for \( x \in L \) with \( \varphi_{I}(x) = [x]_I \), 
        and \( \varphi_{J, I}([x]_{I}) = [x]_{J} \), that \( \varphi_{K, J}([x]_{J}) = \varphi_{K, I}([x]_{I}) \).
        But this holds by Definition, since \( \varphi_{K, J}([x]_{J}) = [x]_{K} \) and \( \varphi_{K, I}([x]_{I}) = [x]_{K} \).
    \end{proof}
    While ideals have a straightfoward relation, which is symmetric to the situation for rings and their quotients, the case of  
    filters in \( L / I \) requires more care.
    For this purpose the following lemma comes in handy:
    \begin{lemma} \label{filterandidealsdonotcontain}
        Let \( L \) be a distributive lattice, for any \( I \in \idl{L} \)
        and \( F \in \flt{L} \), neither \( I \subseteq F \) nor \( F \subseteq I \) holds.
    \end{lemma}
    \begin{proof}
        Suppose the contrary, i.e. \( I \subseteq F \), by Definition
        any \( I \) is downward closed and absorbing with respect to \( \wedge \), while \( F \)
        is upward closed and absorbing with respect to \( \vee \). Take any \( a \in L \)
        and \( i \in I \), if \( a \) is not already in one of them
        then \( a \wedge i \in I \subseteq F \), but \( a \wedge i \leq a \) and \( F \)
        is upward closed, hence \( a \in F \) and by arbitrariness of \( a \), \( F \) can not be proper
        and therefore not a filter.
        The converse for \( F \subseteq \) follows dually with \( a \vee f \) for \( f \in F \).
    \end{proof}
    Note that this does not necessarily mean that all filters and ideals are disjoint! Those who do however
    have a special relation.
    \begin{proposition} \label{filtersinidealquot}
        Let \( I \in \idl{L} \) then there's an inclusion preserving bijection 
        \[ \Psi : \satflt{I}{L} \to \flt{L / I}, \]
        where \( \satflt{I}{L} \) is the set of filters, which are
        disjoint and saturated with respect to I.
        A filter \( F \) is saturated with respect to an ideal \( I \), if
        for \( a \in F \) and \( b \in L \), 
        \[ 
            [a]_{L / I} = [b]_{L / I} \implies b \in F.
        \]
        For such filters, primeness and being an ultrafilter is preserved.
    \end{proposition}
    \begin{proof}
        Let \( \varphi : L \to L / I \) be the canonical projection and define \( \Psi(F) := \varphi(F) \) in the sense of sets.
        \begin{itemize}
            \item \( F \in \satflt{I}{L} \implies \Psi(F) \in \flt{L / I} \).
                \begin{itemize}
                    \item 
                        Take \( F \in \satflt{I}{L} \) with \( a, b \in F \)
                        then \( \varphi(a) \wedge \varphi(b) = \varphi(a \wedge b) \)
                        so for \( [a], [b] \in \Psi(F) \) also \( [a] \wedge [b] = [a \wedge b] \in \Psi(F) \).
                        Similarly for \( a \in F \) and \( b \in L \), we can argue \( [a \vee b] \in \Psi(F) \).
                        Further for \( a \in F \) and \( b \in L \), \( [a] \leq [b] \iff [a \vee b] = [b] \)
                        and \( F \) is a filter, so \( a \vee b \in F \), but then by saturation \( b \in F \), so \( [b] \in \Psi(F) \).
                \end{itemize}
            \item \( F \in \satflt{I}{L} \) is prime, then \( \Psi(F) \in \flt{L / I} \) is as well.
                \begin{itemize}
                    \item 
                        By the previous \( \Psi(F) \in \flt{L / I} \).
                        If \([a \vee b] \in \Psi(F) \) for \( a, b \in L \) then there are \( f \in F \) and \( i \in I \)
                        s.t. \( (a \vee b) \vee i = f \vee i \), \( F \) is absorbing with respect to \( \vee \), so \( f \vee i \in F \)
                        and therefore \( (a \vee b) \vee i \in F \). Now \( F \) is prime and disjoint from \( I \), so only \( (a \vee b) \in F \)
                        is possible and applying primeness again gives \( a \in F \) or \( b \in F \), but then one of their equivalence classes is in \( \Psi(F) \).
                \end{itemize}
            \item \( F' \in \flt{L / I} \implies \Psi^{-1}(F') \in \flt{L} \) and saturated with \( I \).
                \begin{itemize}
                    \item 
                        Let \( F' \in \flt{L / I} \), if \( a, b \in \Psi^{-1}(F') \) then \( [a], [b] \in F' \)
                        and as \( F' \) is a filter, \( [a \wedge b] \in F' \) but then \( a \wedge b \in \Psi^{-1}(F') \).
                        Similarly for \( a \in \Psi^{-1}(F') \) and \( b \in L \), it holds \( [a \vee b] \in F' \),
                        and we conclude \( a \vee b \in \Psi^{-1}(F') \). Combining the results we have that \( \Psi^{-1}(F') \) is a filter.
                        Let \( a \in \Psi^{-1}(F') \) and \( b \in L \), s.t. 
                        \( [a] = [b] \in F' \), trivially then also \( b \in \Psi^{-1}(F') \), making \( \Psi^{-1}(F') \) saturated.
                \end{itemize}
            \item \( F' \in \flt{L / I} \) is prime then \( \Psi^{-1}(F') \in \flt{L} \) is as well.
                \begin{itemize}
                    \item 
                        Suppose that \( F' \in \flt{L / I} \) is also prime, by the previous \( \Psi^{-1}(F') \in \satflt{I}{L}\),
                        assume \( a \vee b \in \Psi^{-1}(F')\), then \( [a \vee b] \in F' \) and due to assumption \( [a], [b] \in F' \)
                        by saturation of \( \Psi^{-1}(F') \) it follows that \( a \in \Psi^{-1}(F') \) or \( b \in \Psi^{-1}(F') \).
                \end{itemize}
            \item \( F' \in \flt{L / I} \implies \Psi^{-1}(F') \in \satflt{I}{L} \).
                \begin{itemize}
                    \item 
                        Assume for contradiction that there is an \( i \in I \cap \Psi^{-1}(F') \), then \( [i] \in F' \) and as
                        all elements in \( I \) are in the same quotient class, also \( [j] \in F' \) for all \( j \in I \setminus \{i\}\),
                        but then \( I \subseteq \Psi^{-1}(F') \), and \( \Psi^{-1}(F') \) can not have been a filter according to Lemma \ref{filterandidealsdonotcontain}
                        and therefore \( I \) and \( \Psi^{-1}(F') \) are disjoint.
                \end{itemize}
        Hence \( \Psi \) is well defined bijective map with the asserted domain and codomain.
        In the proof we defined \( \Psi \) even on the level of sets instead of filters hence it
        is inclusion preserving and therefore ultrafilters in \( \satflt{I}{L} \) are preserved.
        \end{itemize}
    \end{proof}
    \begin{remark} \label{dualfilterquotient}
        All of the previous constructions also exist for filters, which one might argue is an immediate consequence by
        the duality principle (see also the end of Section 1.1 of \cite{genlattheorygraetzer}). 
        If \( F \subset L \) is a filter, then 
        we can construct the quotient lattice \( L / F \) under the equivalence relation 
        \[
            \rho_1 \sim \rho_2 \iff \exists \eta \in F, \rho_1 \wedge \eta = \rho_2 \wedge \eta.
        \]
        Then everything in \( F \) is sent to \( \top \) in \( L / F \). Additionally all filters
        containing \( F \) are sent to filters in \( L / F \) and the properties of being an ultrafilter
        or prime (that is \( a \vee b \in F \iff a \in F \text{ or } b \in F \)) are maintained. 
        Also for filters \( F_1 \subseteq F_2 \) there is a well-defined lattice homomorphism \( \varphi : L / F_2 \to L / F_1 \)
        and disjoint saturated ideals from a filter correspond to ideals in the filter quotient.
    \end{remark}
\end{subsection}
